% Half-title, title and copyright pages.
%
% \bookversion is the release the PDF was built for: `make VERSION=book-v2` passes it in, and the
% release workflow passes the tag. Without it the title page carries only the date of the build.
\providecommand{\bookversion}{}
\newcommand{\bookdate}{\number\day\space\ifcase\month\or January\or February\or March\or April\or
  May\or June\or July\or August\or September\or October\or November\or December\fi\space
  \number\year}
\newcommand{\bookedition}{\ifx\bookversion\empty\else\bookversion\,·\,\fi\bookdate}

\thispagestyle{empty}
\vspace*{2.2in}
\begin{flushright}
  {\Large\bfseries UART: Hardware, Driver \& Integration}
\end{flushright}
\cleardoublepage

\thispagestyle{empty}
\vspace*{1.1in}
\begin{flushleft}
  {\color{accent}\rule{0.9in}{2.5pt}}\par\vspace{1.4ex}
  {\fontsize{32}{38}\selectfont\bfseries UART\par}
  \vspace{0.6ex}
  {\fontsize{24}{29}\selectfont\bfseries Hardware, Driver\\[0.15ex] \& Integration\par}
  \vspace{2.2ex}
  {\Large\itshape FPGA meets MCU: a peripheral in VHDL,\\its driver in C++, and the bench between
    them\par}
  \vspace{1.0in}
  {\large Erik Pihl\par}
  \vfill
  {\small\scshape\lsstyle qrtech academy\par}
  \vspace{0.6ex}
  {\small\color{muted} \bookedition\par}
\end{flushleft}
\newpage

\thispagestyle{empty}
\vspace*{\fill}
{\small
\noindent\textit{UART: Hardware, Driver \& Integration}\\
Copyright \textcopyright\ 2026 QRTECH Academy.

\medskip
\noindent The text and figures of this book, and of the course it is typeset from, are licensed
under the Creative Commons Attribution-NonCommercial-ShareAlike 4.0 International License
(CC BY-NC-SA 4.0). You may copy and redistribute them, and adapt them, for any non-commercial
purpose, provided you give appropriate credit and share what you make under the same license. To
view a copy of the license, visit \url{https://creativecommons.org/licenses/by-nc-sa/4.0/}.

\medskip
\noindent The code in this book may also be used under the MIT License, which is the license of all
the code in the companion repository. The repository holds the lecture material, the protocol
specification, the provided VHDL and its testbenches, the host test suites and the build scripts:

\medskip
\noindent\url{https://github.com/qrtech-academy/uart}

\medskip
\noindent The VHDL in this book is VHDL-93, simulated with GHDL and synthesized with Quartus Prime
Lite. The C++ is C++17 on the host and, wherever it runs on the ATmega328P, a subset that also
compiles as C++14 with avr-gcc. Arduino Nano, ATmega328P, DE0-CV, Cyclone~V, Quartus, USB-Blaster,
CP2102 and FT232 are trademarks of their respective owners.

\medskip
\noindent Set in TeX Gyre Pagella and DejaVu Sans Mono with Lua\LaTeX.

\medskip
\noindent\textcolor{muted}{Typeset from the ten lectures of the course, their exercises, the UART
register protocol, and the two self-assessment papers. This edition: \bookedition.}
\par}
\cleardoublepage
